\documentclass[10pt,a4paper]{amsart}
\usepackage[margin=19mm]{geometry}
\usepackage{amsmath,amssymb,mathtools}
\usepackage[scr=rsfs,cal=euler]{mathalfa}
\usepackage{xfrac}
\usepackage[unicode,hidelinks,bookmarks=false]{hyperref}
\newcommand{\ie}{i.e.\ }
\newcommand{\cf}{cf.\ }
\newcommand{\resp}{respectively\ }
\newcommand{\Subsection}{\subsection}
\providecommand{\nobreakdash}{\nobreak\hbox{-}}
\setlength{\emergencystretch}{2em}
\multlinegap0pt
\usepackage{bm}
\usepackage{bbm}
\usepackage{dsfont}
\usepackage{xspace}
\usepackage{cancel}
\usepackage{mathtools}
\usepackage{amsthm}
\makeatletter
\@ifundefined{lemma}{\newtheorem{lemma}{Lemma}}{}
\@ifundefined{remark}{\newtheorem{remark}{Remark}}{}
\makeatother
\makeatletter
\newcommand{\Hon}[1]{\|#1\|_1}
\newcommand{\Ltw}[1]{\|#1\|_0}
\newcommand{\LtwD}[2]{\|#1\|_{0,#2}}
\newcommand{\LtwE}[1]{\|#1\|_{0,E}}
\newcommand{\LtwT}[1]{\|#1\|_{0,T}}
\newcommand{\anorm}[1]{\|#1\|_a}
\@ifundefined{breakablealgorithm}{\newenvironment{breakablealgorithm}
{% \begin{breakablealgorithm}
		\begin{center}
			\refstepcounter{algorithm}% New algorithm
			\hrule height.8pt depth0pt \kern2pt% \@fs@pre for \@fs@ruled
			\renewcommand{\caption}[2][\relax]{% Make a new \caption
				{\raggedright\textbf{\ALG@name~\thealgorithm} ##2\par}%
				\ifx\relax##1\relax % #1 is \relax
				\addcontentsline{loa}{algorithm}{\protect\numberline{\thealgorithm}##2}%
				\else % #1 is not \relax
				\addcontentsline{loa}{algorithm}{\protect\numberline{\thealgorithm}##1}%
				\fi
				\kern2pt\hrule\kern2pt
			}
		}{% \end{breakablealgorithm}
		\kern2pt\hrule\relax% \@fs@post for \@fs@ruled
	\end{center}
}}{}
\newcommand{\est}{\mathrm{est}}
\def\geq{\geqslant}
\newcommand{\ha}{\hat{a}}
\newcommand{\hmcL}{\hat{\mathcal{L}}}
\newcommand{\hu}{\hat{u}}
\def\leq{\leqslant}
\newcommand{\mbV}{\mathbb{V}}
\newcommand{\mcL}{\mathcal{L}}
\newcommand{\mcN}{\mathcal{N}}
\newcommand{\mcT}{\mathcal{T}}
\newcommand{\reduce}{\mathrm{red}}
\newcommand{\rem}{\mathrm{rem}}
\newcommand{\tu}{\tilde{u}}
\newcommand{\valpha}{\boldsymbol{\alpha}}
\newcommand{\vbeta}{\boldsymbol{\beta}}
\newcommand{\vect}[1]{\boldsymbol{#1}}
\newcommand{\veps}{\varepsilon}
\newcommand{\vn}{\boldsymbol{n}}
\makeatother
\title[A correction adaptive two-grid finite element method for non]{A correction adaptive two-grid finite element method for nonselfadjoint or indefinite elliptic problems}
\author{Fei Li, Qingguo Hong, Ming Tang, Liuqiang Zhong}
\date{Source: arXiv:2604.24567v1 (CC BY 4.0)}
\begin{document}
\maketitle

\noindent\emph{Source and licence.} A correction adaptive two-grid finite element method for nonselfadjoint or indefinite elliptic problems. Version arXiv:2604.24567v1 (CC BY 4.0), distributed under the Creative Commons Attribution 4.0 International licence, as recorded in the arXiv licence metadata for this exact version and shown on the abstract page https://arxiv.org/abs/2604.24567v1. Full rights notice: licenses/arxiv-2604.24567-CC-BY-4.0.txt.\par\smallskip

\noindent\textit{Editorial scope.} Counted unit: the Lemma whose source label is \texttt{lem-est-reduce} in the article (source0.tex lines 717--727, source label \texttt{lem-est-reduce}), delivered together with the article's own complete original proof of it (source0.tex lines 729--794, one \texttt{proof} environment of 66 lines). No mathematics is written, repaired or re-derived: statement and proof are the article's own text line for line. Required context retained uncounted: model-VP (lines 281--283); equivalence-norm (lines 297--299); solve-SAFEM (lines 352--354); solve-ATGFEM (lines 363--365); mysolver-k-1 (lines 387--389); mysolver-k-2 (lines 391--393); rem-CATG-1 (lines 398--403); def-residual (lines 432--437); jump (lines 439--441); eta-T (lines 446--449); lem-L2-lifting (lines 503--512); result-L2-lifting (lines 505--507). Every definition, display and label of those lines is retained unchanged. Omitted from this package and not redistributed: 1--280, 284--296, 300--351, 355--362, 366--386, 390--390, 394--397, 404--431, 438--438, 442--445, 450--502, 508--716, 728--728, 795--1079. Nothing inside a retained excerpt is omitted or altered: every retained body is delivered exactly as the article's lines are written, so no omission receipt of any kind accompanies this package. Citations: the retained excerpts carry no citation command at all, so no bibliography entry of the article is transcribed and no reference is redistributed. Reference adaptation: none was needed and none was made; every reference inside the delivered unit resolves inside it, so no unresolved reference or citation is produced. Printed slips kept literal: no printed word, symbol, hypothesis or number of the counted statement or of its proof was altered. Layout and class: the article's own class is \texttt{elsarticle} with packages including lineno, stmaryrd, amsmath, amssymb, amsthm, graphicx, epic, eepic, epsfig, color, subfigure, placeins, multirow, varwidth; the wrapper is the lane's pinned \texttt{amsart} editorial wrapper at 10pt on A4, which restates the theorem-like environments the retained text needs and adds the article's own macro definitions that the retained text needs (\texttt{\textbackslash Hon}, \texttt{\textbackslash Ltw}, \texttt{\textbackslash LtwD}, \texttt{\textbackslash LtwE}, \texttt{\textbackslash LtwT}, \texttt{\textbackslash anorm}, \texttt{\textbackslash est}, \texttt{\textbackslash geq}, \texttt{\textbackslash ha}, \texttt{\textbackslash hmcL}, \texttt{\textbackslash hu}, \texttt{\textbackslash leq}, \texttt{\textbackslash mbV}, \texttt{\textbackslash mcL}, \texttt{\textbackslash mcN}, \texttt{\textbackslash mcT}, \texttt{\textbackslash reduce}, \texttt{\textbackslash rem}, \texttt{\textbackslash tu}, \texttt{\textbackslash valpha}, \texttt{\textbackslash vbeta}, \texttt{\textbackslash vect}, \texttt{\textbackslash veps}, \texttt{\textbackslash vn}), with extra wrapper packages (bm, bbm, dsfont, xspace, cancel, mathtools). The delivered numbering is the unit's own. Assets: no retained span contains a figure environment, a graphics-inclusion command, a PDF-page inclusion, a table or a TikZ picture; no image file is copied into this package, the package declares no asset dependency and no image file is redistributed with it. Licence: version arXiv:2604.24567v1 of this article is distributed under the Creative Commons Attribution 4.0 International licence, as recorded in the arXiv licence metadata for this exact version and shown on the abstract page https://arxiv.org/abs/2604.24567v1. A full rights notice is delivered at licenses/arxiv-2604.24567-CC-BY-4.0.txt. Characters: every retained excerpt is pure ASCII, so the delivered engine, XeLaTeX with its default Latin Modern text font, reports no missing character.
\begin{equation}\label{model-VP}
	\ha(u,v) = \langle f , v \rangle, \quad \forall\; v \in H_0^1(\Omega), 
\end{equation}

\begin{equation}\label{equivalence-norm}
	 \anorm{v} \simeq \Hon{v}. 
\end{equation}

	\begin{equation}\label{solve-SAFEM}
		\ha(\hu_k,v_k) = \langle f,v_k \rangle, \quad \forall \; v_k \in \mbV_k. 
	\end{equation}

	\begin{equation}\label{solve-ATGFEM}
		a(\tu_k,v_k) = \langle f,v_k \rangle - N(\tu_{k-1},v_k), \quad \forall \; v_k \in \mbV_k, 
	\end{equation}

	\begin{equation}\label{mysolver-k-1}
		\ha(e_{0,k-1},v_0) = \langle f,v_0 \rangle - \ha(u_{k-1},v_0), \quad \forall \; v_0 \in \mbV_0.
	\end{equation}

	\begin{equation}\label{mysolver-k-2}
		a(u_k,v_k) = (f,v_k) - N(u_{k-1}+e_{0,k-1},v_k), \quad \forall\; v_k \in \mbV_k. 
	\end{equation}

\begin{remark}\label{rem-CATG-1}For the case of $k = 1$, the right-hand side of \eqref{mysolver-k-1} is zero due to $u_0 = \hu_0$ being the solution of \eqref{solve-SAFEM}. This implies that $e_{0,0} = 0$, which means that \eqref{mysolver-k-1} does not work, and \eqref{mysolver-k-2} will be transformed into 
\begin{equation}\label{solve-CATGFEM-1}
	a(u_1,v_1) = \langle f, v_1 \rangle - N(u_0,v_1), \quad \forall \; v_1 \in \mbV_1, 
\end{equation}
which is indeed \eqref{solve-ATGFEM} for $k=1$ of the ATGFEM due to $u_0 = \tu_0$.  
\end{remark}

\begin{equation}\label{def-residual}
	R_T(u_k) = \begin{cases}
		f|_T - \hmcL u_0 |_T, & k = 0, \\
		f|_T - \mcL u_k|_T - \mcN u_{k-1}^* |_T, & k \geq 1, 
	\end{cases}
\end{equation}

\begin{equation}\label{jump}
	J_E(u_k) := [\valpha \nabla u_k \cdot \vect{n}]|_E := \valpha\nabla u_k|_{T_1} \cdot \vect{n}_1 + \valpha \nabla u_k |_{T_2} \cdot \vect{n}_2,
\end{equation}

\begin{equation}\label{eta-T}
	\eta_k(u_k,T) := \Big(h_T^2 \|R_T(u_k)\|_{0,T}^2 + h_T 
	\sum_{E \subset \partial T \cap \Omega}\LtwE{J_E(u_k)}^2\Big)^\frac{1}{2}, 
\end{equation}	

\begin{lemma}[$L^2$-lifting property]\label{lem-L2-lifting}
	Let $u$, $e_{0,k-1}$ and $u_k$ be the solutions of \eqref{model-VP}, \eqref{mysolver-k-1} and \eqref{mysolver-k-2}, and keep $u_{k-1}^*=e_{0,k-1} + u_{k-1}$ in mind. Then, there holds 
		\begin{equation}\label{result-L2-lifting}
			\Ltw{u-u_{k-1}^*} \leq \kappa(h_0) \anorm{u-u_{k-1}}
		\end{equation}
with a positive constant $\kappa(h_0) = \kappa(h_0,\Omega,\valpha,\vbeta,\gamma)$ satisfying
		\begin{equation}\label{property-kappaH}
			\lim_{h_0 \rightarrow 0} \kappa(h_0) = 0.  
		\end{equation}
\end{lemma}

\begin{lemma}[Estimator reduction]\label{lem-est-reduce} 
	Let $u$ and $u_k$ be the solutions of \eqref{model-VP} and \eqref{mysolver-k-2}. Then, for $k \geq 2$, there holds
	\begin{align}
		\eta_k^2(u_k) \leqslant~& \rho_\est \eta_{k-1}^2(u_{k-1}) 
+ C_\est\anorm{u_k-u_{k-1}}^2+ C_\rem \kappa(h_0)^2\left(\anorm{u-u_{k-1}}^2 + \anorm{u-u_{k-2}}^2\right)
\label{result-est-reduce}
	\end{align}
with $\rho_\est = \rho_\est(d,\theta) \in (0,1)$, and the positive constants $C_\est = C_\est(d,\Omega,\valpha)$ and $C_\rem =  C_\rem(\vbeta,\gamma)$ both also dependent of the uniform shape regularity of the mesh family $\{\mcT_k\}_{k \geq 0}$. 

Particularly for $k = 1$, the term $C_\rem \kappa(h_0)^2(\anorm{u-u_{k-1}}^2 + \anorm{u-u_{k-2}}^2)$ of \eqref{result-est-reduce} will vanish. 
\end{lemma}

\begin{proof}For the elementwise error estimator \eqref{eta-T}, the following inequality 
	$$
	(b_1^{2}+c_1^{2})^{\frac{1}{2}}-(b_2^{2}+c_2^{2})^{\frac{1}{2}}\leqslant (b_1-b_2) + (c_1 - c_2), \quad \forall \; b_1,b_2,c_1,c_2 \geqslant 0, 
	$$
implies that 
	\begin{align}
	\eta_k(u_k,T) - \eta_k(u_{k-1},T)
		\leq~&h_T\LtwT{R_T(u_k)-R_T(u_{k-1})}+ \sum_{E \subset \partial T\cap\Omega}h_T^\frac{1}{2}\LtwE{J_E(u_k)-J_E(u_{k-1})}. \label{step1-stability}
	\end{align}
By the definition \eqref{def-residual} of $R_T(\cdot)$ and the inverse estimate, the first term in the right-hand side of \eqref{step1-stability} is bounded by 
	\begin{align}
	h_T \LtwT{R_T(u_k)-R_T(u_{k-1})}  \leq~&
		h_T\Big(\LtwT{(\nabla\cdot\valpha)\cdot\nabla(u_k-u_{k-1})}
		+\LtwT{\valpha:\nabla^2(u_k-u_{k-1})}
		\nonumber\\
		&+\LtwT{\vbeta\cdot\nabla(u_{k-1}^*-u_{k-2}^*)}
		 +\LtwT{\gamma(u_{k-1}^*-u_{k-2}^*)}\Big)
		\nonumber\\
		\lesssim~& \LtwT{\nabla(u_k-u_{k-1})} + \LtwT{u_{k-1}^* - u_{k-2}^*}
		\label{step2-stability}
	\end{align}
for $k \geq 2$, where $\nabla^2(\cdot)$ denotes the Hessian matrix. Note that, for $k = 1$, the last term in the right-hand side of \eqref{step2-stability} will vanish due to $u_0^* = u_0 + e_{0,0} = u_0$ from Remark \ref{rem-CATG-1}. 

For the estimation of the second term in the right-hand side of \eqref{step1-stability}, we note that for any side $E \in \partial T \cap \Omega$, there exists an element $T' \in \mcT_k$ such that $T' \cap T=E$. By the definition \eqref{jump} of the jump, the trace theorem and the inverse estimate, we show that 
\begin{align}
	h_T^{\frac{1}{2}}\LtwE{J_E(u_k)-J_E(u_{k-1})}\leq~&
		h_T^\frac12(\LtwE{\valpha\nabla(u_k-u_{k-1})|_{\partial T \cap E}\cdot\vn_T}+\LtwE{\valpha\nabla(u_k-u_{k-1})|_{\partial T' \cap E}\cdot\vn_{T'}})
			\nonumber\\
		\lesssim~& \LtwD{\nabla(u_k-u_{k-1})}{T \cup T'}. 
			\nonumber
\end{align}
Then, it holds that 
	\begin{align}
		\sum_{E \in \partial T \cap \Omega}h_T^\frac12\LtwE{J_E(u_k)-J_E(u_{k-1})}
		\lesssim~&(d+1)\LtwD{\nabla(u_k-u_{k-1})}{\omega_T}, 
		\label{step3-stability}
	\end{align}
where $\omega_T$ denotes the set of elements sharing the edge $E \in \partial T \cap \Omega$ with the element $T$.  
Substituting \eqref{step2-stability} and \eqref{step3-stability} into \eqref{step1-stability} yields that 
	\begin{align}
		\eta_k(u_k,T) \leq~&  \eta_k(u_{k-1},T) +  C_1\LtwD{\nabla (u_k-u_{k-1})}{\omega_T} + C_2\LtwT{u_{k-1}^*-u_{k-2}^*} \nonumber
	\end{align}
for $k \geq 2$ with positive constants $C_1= C_1(d,\Omega,\valpha)$ and $C_2 = C_2(\vbeta,\gamma)$ both also dependent of the uniform shape regularity of the mesh family $\{\mcT_k\}_{k \geq 0}$, where the last term in the above right-hand side will vanish for $k = 1$. 

Applying the Young inequality with $\veps > 0$ to the square of the above inequality, summing the arising result over all $T \in \mcT_k$ and using the norm equivalence \eqref{equivalence-norm} as well as \eqref{result-L2-lifting} of Lemma \ref{lem-L2-lifting} lead to that 
\begin{align}
	\eta_k^2( u_k) \leq~& (1+\veps)\eta_k^2(u_{k-1})
		+2(1+\veps^{-1})\Big((d+2)C_1^2\Ltw{\nabla(u_k-u_{k-1})}^2 
		+C_2^2\Ltw{u_{k-1}^*-u_{k-2}^*}^2\Big)\nonumber\\
		\leq~&(1+\veps)\eta_k^2(u_{k-1}) + 2(1+\veps^{-1})(d+2)C_1^2C(\Omega,\valpha)
		\anorm{u_k-u_{k-1}}^2 \nonumber\\
		&+4(1+\veps^{-1})C_2^2\kappa(h_0)^2(\anorm{u-u_{k-1}}^2 + \anorm{u-u_{k-2}}^2)
		\label{step4-stability} 
\end{align}
for $k \geq 2$, where we use the fact that the sum of norms on $\omega_T$ covers each element at least $d+2$ times. While the term $4(1+\veps^{-1})C_2^2\kappa(h_0)^2(\anorm{u-u_{k-1}}^2 + \anorm{u-u_{k-2}}^2)$ in \eqref{step4-stability} will vanish for $k=1$. 

 Furthermore, it is easy to obtain that 
 \begin{equation}\label{reduction}
 	\eta_k^2(u_{k-1}) \leq \rho_\reduce\eta_{k-1}^2(u_{k-1}), 
 \end{equation}
 with the constant $\rho_\reduce:= 1 - 2^{-1/d}\theta$. Subsituting \eqref{reduction} into \eqref{step4-stability} and taking $\veps$ sufficiently small such that 
 $\rho_\est := (1+\veps)\rho_\reduce \in (0,1)$ conclude the proof. 
% 
% with the constant $\rho_\reduce=1-(1-2^{-\frac{1}{d}})\theta \in (0,1)$ due to the marking parameter $\theta \in (0,1)$. Substituting \eqref{reduction} into \eqref{step4-stability} and letting $\veps_1$ be sufficiently small such that 
% $(1+\veps)\rho_\reduce \in (0,1)$ concludes the proof. 
\end{proof}

\end{document}
